\chapter{Resultados}\label{cap4-resultados}

Este capítulo apresenta os resultados experimentais obtidos a partir da execução das quatro variantes do algoritmo memético AG+\(2\)-opt (CPU, híbrida ingênua, híbrida otimizada e totalmente em GPU) sobre o conjunto completo de \totalinstancias\ instâncias da TSPLIB.
A análise é dividida em quatro partes principais:

\begin{enumerate}
  \item
        uma avaliação detalhada por instância, destacando a progressão entre \(\menorinstancia\) e \(\maiorinstancia\) cidades (com a tabela integral das \totalinstancias\ instâncias no \apref{apendice-resultados});
  \item
        uma análise agregada sobre todo o conjunto avaliado, focada nas métricas de performance e qualidade das soluções;
  \item
        a aplicação dos testes estatísticos não paramétricos descritos no \capref{cap3-materiais-metodos}, acompanhada dos diagramas de diferença crítica (tempo e \emph{gap}) e da análise de escalonamento; e
  \item
        a divisão do tempo de execução entre as oito etapas do ciclo evolutivo e a análise de gargalos com base na Lei de Amdahl.
\end{enumerate}

\section{Resultados experimentais}\label{sec-41-resultados-experimentais}

Os experimentos foram conduzidos conforme o protocolo definido na \secref{sec-33-selecao-instancias}, com \totalrepeticoes\ repetições independentes pareadas (sementes \(\sementeinicial\) a \(\sementefinal\)) para cada par algoritmo--instância, totalizando \(\totalcombinacoes\) combinações e \(\totalexecucoes\) execuções registradas.
Graças à busca local vetorizada associada à lista de candidatos \(k_{\text{viz}} = \kviz\), todas as quatro variantes foram executadas integralmente sobre todas as \totalinstancias\ instâncias selecionadas.

As métricas primárias coletadas foram:

\begin{itemize}
  \item
        \textbf{Tempo Médio de Execução (s):} Tempo total para completar a execução por repetição.
  \item
        \textbf{Gap Médio (\%):} Desvio percentual do custo da melhor solução encontrada em relação ao ótimo conhecido da instância.
  \item
        \textbf{Performance:} Razão entre o tempo de execução da variante de referência (CPU) e a variante avaliada.
\end{itemize}

\section{Análise por instância}\label{sec-42-analise-por-instancia}

A \tabref{tbl:results-detailed} reúne os resultados detalhados de sete instâncias selecionadas entre \(52\) e \(\maiorinstancia\) cidades (\texttt{berlin52}, \texttt{kroA100}, \texttt{d198}, \texttt{pr264}, \texttt{att532}, \texttt{si535} e \texttt{pr1002}); os valores completos de todas as \totalinstancias\ instâncias (de \(\menorinstancia\) a \(\maiorinstancia\) cidades) constam na \tabref{tbl:results-all-instances} do \apref{apendice-resultados}.

\begin{tabelapadrao}
  \begin{longtable}[]{@{}
      >{\raggedright\arraybackslash}p{(\linewidth - 12\tabcolsep) * \real{0.11}}
      >{\centering\arraybackslash}p{(\linewidth - 12\tabcolsep) * \real{0.07}}
      >{\raggedright\arraybackslash}p{(\linewidth - 12\tabcolsep) * \real{0.24}}
      >{\centering\arraybackslash}p{(\linewidth - 12\tabcolsep) * \real{0.26}}
      >{\centering\arraybackslash}p{(\linewidth - 12\tabcolsep) * \real{0.10}}
      >{\centering\arraybackslash}p{(\linewidth - 12\tabcolsep) * \real{0.11}}
      >{\centering\arraybackslash}p{(\linewidth - 12\tabcolsep) * \real{0.11}}@{}}
    \caption{Resultados detalhados por instância (seleção representativa)}\label{tbl:results-detailed}                                                   \\
    \toprule
    Instância         & \(n_{coords}\) & Algoritmo         & Custo médio (\(\pm \sigma\))  & \emph{Gap} médio (\%) & Tempo médio (s) & Ganho de perf.    \\
    \midrule
    \endfirsthead
    \toprule
    Instância         & \(n_{coords}\) & Algoritmo         & Custo médio (\(\pm \sigma\))  & \emph{Gap} médio (\%) & Tempo médio (s) & Ganho de perf.    \\
    \midrule
    \endhead
    \bottomrule
    \endlastfoot
    \textbf{berlin52} & 52             & CPU               & \(7553{,}03 \pm 17{,}39\)     & \(0{,}15\)            & \(1{,}59\)      & \(1{,}00\times\)  \\*
                      & 52             & Híbrida Ingênua   & \(7545{,}23 \pm 12{,}25\)     & \(0{,}04\)            & \(0{,}51\)      & \(3{,}12\times\)  \\*
                      & 52             & Híbrida Otimizada & \(7545{,}23 \pm 12{,}25\)     & \(0{,}04\)            & \(0{,}08\)      & \(20{,}23\times\) \\*
                      & 52             & Totalmente em GPU & \(7550{,}97 \pm 19{,}97\)     & \(0{,}12\)            & \(0{,}02\)      & \(75{,}03\times\) \\
    \midrule
    \textbf{kroA100}  & 100            & CPU               & \(21389{,}43 \pm 54{,}23\)    & \(0{,}50\)            & \(4{,}81\)      & \(1{,}00\times\)  \\*
                      & 100            & Híbrida Ingênua   & \(21391{,}40 \pm 63{,}09\)    & \(0{,}51\)            & \(1{,}18\)      & \(4{,}06\times\)  \\*
                      & 100            & Híbrida Otimizada & \(21391{,}40 \pm 63{,}09\)    & \(0{,}51\)            & \(0{,}19\)      & \(25{,}73\times\) \\*
                      & 100            & Totalmente em GPU & \(21372{,}60 \pm 56{,}96\)    & \(0{,}43\)            & \(0{,}05\)      & \(94{,}86\times\) \\
    \midrule
    \textbf{d198}     & 198            & CPU               & \(15918{,}27 \pm 16{,}32\)    & \(0{,}88\)            & \(14{,}66\)     & \(1{,}00\times\)  \\*
                      & 198            & Híbrida Ingênua   & \(15916{,}07 \pm 14{,}70\)    & \(0{,}86\)            & \(3{,}08\)      & \(4{,}76\times\)  \\*
                      & 198            & Híbrida Otimizada & \(15916{,}07 \pm 14{,}70\)    & \(0{,}86\)            & \(0{,}55\)      & \(26{,}84\times\) \\*
                      & 198            & Totalmente em GPU & \(15912{,}93 \pm 18{,}04\)    & \(0{,}84\)            & \(0{,}17\)      & \(88{,}62\times\) \\
    \midrule
    \textbf{pr264}    & 264            & CPU               & \(49552{,}43 \pm 71{,}20\)    & \(0{,}85\)            & \(23{,}62\)     & \(1{,}00\times\)  \\*
                      & 264            & Híbrida Ingênua   & \(49551{,}47 \pm 98{,}13\)    & \(0{,}85\)            & \(4{,}81\)      & \(4{,}91\times\)  \\*
                      & 264            & Híbrida Otimizada & \(49551{,}47 \pm 98{,}13\)    & \(0{,}85\)            & \(1{,}02\)      & \(23{,}15\times\) \\*
                      & 264            & Totalmente em GPU & \(49527{,}90 \pm 80{,}20\)    & \(0{,}80\)            & \(0{,}34\)      & \(69{,}92\times\) \\
    \midrule
    \textbf{att532}   & 532            & CPU               & \(28000{,}10 \pm 83{,}69\)    & \(1{,}13\)            & \(158{,}24\)    & \(1{,}00\times\)  \\*
                      & 532            & Híbrida Ingênua   & \(27981{,}53 \pm 57{,}22\)    & \(1{,}07\)            & \(29{,}33\)     & \(5{,}39\times\)  \\*
                      & 532            & Híbrida Otimizada & \(27981{,}53 \pm 57{,}22\)    & \(1{,}07\)            & \(6{,}67\)      & \(23{,}73\times\) \\*
                      & 532            & Totalmente em GPU & \(27999{,}93 \pm 78{,}55\)    & \(1{,}13\)            & \(3{,}39\)      & \(46{,}63\times\) \\
    \midrule
    \textbf{si535}    & 535            & CPU               & \(48903{,}20 \pm 21{,}15\)    & \(0{,}94\)            & \(84{,}19\)     & \(1{,}00\times\)  \\*
                      & 535            & Híbrida Ingênua   & \(48900{,}00 \pm 23{,}95\)    & \(0{,}93\)            & \(10{,}97\)     & \(7{,}68\times\)  \\*
                      & 535            & Híbrida Otimizada & \(48900{,}00 \pm 23{,}95\)    & \(0{,}93\)            & \(2{,}48\)      & \(33{,}99\times\) \\*
                      & 535            & Totalmente em GPU & \(48909{,}00 \pm 18{,}43\)    & \(0{,}95\)            & \(1{,}23\)      & \(68{,}23\times\) \\
    \midrule
    \textbf{pr1002}   & 1002           & CPU               & \(263558{,}83 \pm 771{,}47\)  & \(1{,}74\)            & \(682{,}58\)    & \(1{,}00\times\)  \\*
                      & 1002           & Híbrida Ingênua   & \(263288{,}90 \pm 954{,}07\)  & \(1{,}64\)            & \(115{,}35\)    & \(5{,}92\times\)  \\*
                      & 1002           & Híbrida Otimizada & \(263288{,}90 \pm 954{,}07\)  & \(1{,}64\)            & \(25{,}55\)     & \(26{,}72\times\) \\*
                      & 1002           & Totalmente em GPU & \(263409{,}57 \pm 1121{,}29\) & \(1{,}68\)            & \(14{,}03\)     & \(48{,}66\times\) \\
  \end{longtable}
  \notatabela{Médias sobre as \totalrepeticoes\ repetições. Os ganhos de velocidade (\textit{Ganho de perf.}) foram calculados a partir dos tempos médios exatos antes do arredondamento a duas casas decimais (por exemplo, em \texttt{berlin52}, \(1{,}594134\,\text{s} / 0{,}021246\,\text{s} \approx 75{,}03\times\)).}
\end{tabelapadrao}

Nas instâncias menores e intermediárias da tabela (\texttt{berlin52}, \texttt{kroA100} e \texttt{d198}), a variante \textbf{Totalmente em GPU} atinge os maiores ganhos relativos sobre a CPU: \(94{,}86\times\) em \texttt{kroA100} (\(0{,}05\)\,s contra \(4{,}81\)\,s), \(88{,}62\times\) em \texttt{d198} (\(0{,}17\)\,s contra \(14{,}66\)\,s) e \(75{,}03\times\) em \texttt{berlin52} (\(0{,}02\)\,s contra \(1{,}59\)\,s) --- chegando a \(112{,}48\times\) em \texttt{gr96} na tabela integral do \apref{apendice-resultados}.
A \textbf{Híbrida Otimizada} ocupa consistentemente o segundo lugar (\(20{,}23\times\) a \(26{,}84\times\) nessas três instâncias), muito à frente da \textbf{Híbrida Ingênua} (\(3{,}12\times\) a \(4{,}76\times\)).

Em \texttt{pr264} (\(264\) cidades), a mesma hierarquia se mantém: a \textbf{Totalmente em GPU} conclui em \(0{,}34\)\,s (\(69{,}92\times\) mais rápida que a CPU), seguida pela \textbf{Híbrida Otimizada} (\(1{,}02\)\,s; \(23{,}15\times\)) e pela \textbf{Híbrida Ingênua} (\(4{,}81\)\,s; \(4{,}91\times\)).
Avançando para \texttt{att532} (\(532\)), \texttt{si535} (\(535\)) e \texttt{pr1002} (\(\maiorinstancia\)), o ganho em tempo absoluto fica ainda mais evidente: em \texttt{pr1002}, a \textbf{Totalmente em GPU} reduz o tempo médio de \(682{,}58\)\,s na CPU para \(14{,}03\)\,s (\(48{,}66\times\) mais rápida que a CPU, \(8{,}22\times\) que a Híbrida Ingênua com \(115{,}35\)\,s e \(1{,}82\times\) que a Híbrida Otimizada com \(25{,}55\)\,s).
Em relação à qualidade das soluções, os \emph{gaps} permanecem praticamente iguais entre as variantes ao longo de toda a tabela --- em \texttt{pr1002}, por exemplo, eles ficam em \(1{,}74\%\) na CPU, \(1{,}64\%\) nas duas variantes híbridas (idênticas bit a bit nas \totalrepeticoes\ sementes) e \(1{,}68\%\) na Totalmente em GPU.

\section{Análise agregada, testes estatísticos e escalonamento}\label{sec-43-analise-agregada}

Ao agregar os resultados de todas as \totalinstancias\ instâncias, emergem padrões claros de desempenho que caracterizam o comportamento temporal e a qualidade das soluções.

\begin{tabelapadrao}
  \begin{longtable}[]{@{}
      >{\raggedright\arraybackslash}p{(\linewidth - 6\tabcolsep) * \real{0.25}}
      >{\centering\arraybackslash}p{(\linewidth - 6\tabcolsep) * \real{0.31}}
      >{\centering\arraybackslash}p{(\linewidth - 6\tabcolsep) * \real{0.22}}
      >{\centering\arraybackslash}p{(\linewidth - 6\tabcolsep) * \real{0.22}}@{}}
    \caption{Desempenho agregado sobre todas as \totalinstancias\ instâncias}\label{tbl:results-aggregated}           \\
    \toprule
    Algoritmo         & Tempo mediano \([Q_{25};\, Q_{75}]\) (s) & \emph{Gap} médio (\%) & Ganho rel. CPU (geom.)     \\
    \midrule
    \endfirsthead
    \toprule
    Algoritmo         & Tempo mediano \([Q_{25};\, Q_{75}]\) (s) & \emph{Gap} médio (\%) & Ganho rel. CPU (geom.)     \\
    \midrule
    \endhead
    \bottomrule
    \endlastfoot
    CPU               & \(23{,}62\;[4{,}81;\, 130{,}73]\)        & \(0{,}90 \pm 0{,}41\) & \(1{,}00\times\)           \\
    Híbrida Ingênua   & \(4{,}81\;[1{,}17;\, 25{,}76]\)          & \(0{,}90 \pm 0{,}44\) & \(4{,}57\times\)           \\
    Híbrida Otimizada & \(1{,}02\;[0{,}18;\, 6{,}02]\)           & \(0{,}90 \pm 0{,}44\) & \(23{,}86\times\)          \\
    Totalmente em GPU & \(\mathbf{0{,}34\;[0{,}05;\, 2{,}72]}\)  & \(0{,}91 \pm 0{,}46\) & \(\mathbf{71{,}03\times}\) \\
  \end{longtable}
  \notatabela{Como os tempos de execução variam em ordens de grandeza entre \(n_{coords} = \menorinstancia\) e \(\maiorinstancia\), o tempo agregado é resumido por mediana e intervalo interquartil (\([Q_{25};\, Q_{75}]\)), evitando a distorção de médias aritméticas sobre amostras fortemente assimétricas; o ganho relativo à CPU utiliza a média geométrica das \totalinstancias\ razões por instância.}
\end{tabelapadrao}

A \tabref{tbl:results-aggregated} resume duas conclusões centrais sobre o comportamento empírico do algoritmo memético:

\begin{enumerate}
  \item
        \textbf{A variante Totalmente em GPU obtém o menor tempo de execução:} Com ganho médio geométrico de \(71{,}03\times\) sobre a CPU, \(15{,}53\times\) sobre a Híbrida Ingênua e \(2{,}98\times\) sobre a Híbrida Otimizada (tempo mediano de \(0{,}34\)\,s com IIQ \([0{,}05;\, 2{,}72]\)\,s contra \(23{,}62\)\,s e IIQ \([4{,}81;\, 130{,}73]\)\,s na CPU), a versão residente em GPU confirma que manter a população, os operadores evolutivos e a busca local na VRAM traz três benefícios combinados: elimina o tráfego de matrizes de rotas pelo barramento PCIe a cada geração, evita lançamentos individuais de \emph{kernels} e remove o gargalo dos operadores executados em CPU (sobretudo o cruzamento, detalhado na \secref{sec-44-decomposicao-fases}).
        A Híbrida Otimizada fica em segundo lugar, com ganho geométrico de \(23{,}86\times\) sobre a CPU e \(5{,}22\times\) sobre a Híbrida Ingênua (mediana de \(1{,}02\)\,s contra \(4{,}81\)\,s da Híbrida Ingênua, que atinge \(4{,}57\times\) sobre a CPU).
  \item
        \textbf{Qualidade das soluções preservada entre arquiteturas:} Com os mesmos parâmetros de busca local (\(\kappa = \maxpassadas\) varreduras e \(k_{\text{viz}} = \kviz\) candidatos), as quatro variantes atingem soluções com qualidade muito próxima: os \emph{gaps} médios agregados ficam em \(0{,}90\% \pm 0{,}41\%\) na CPU, \(0{,}90\% \pm 0{,}44\%\) nas duas variantes híbridas e \(0{,}91\% \pm 0{,}46\%\) na Totalmente em GPU.
\end{enumerate}

Para confirmar estatisticamente essas duas conclusões, aplicaram-se o Eixo~1 (\secref{sec-341-eixo1}) e o Eixo~2 (\secref{sec-342-eixo2}) do protocolo de \textcite{demsar2006statistical} ao nível \(\alpha = 0{,}05\):

\begin{itemize}
  \item
        \textbf{Qualidade das soluções (\emph{gap} percentual):} Não há diferença estatisticamente significativa entre as quatro variantes em nenhum dos dois eixos. No Eixo~2 global, o teste de Friedman com correção de Iman--Davenport não rejeita \(H_0\) (\(\chi_F^2 = 2{,}69\); \(F_F(3, 42) = 0{,}8897\); \(p = 0{,}4543 > 0{,}05\)), os postos médios ficam muito próximos (\(2{,}30\) nas duas híbridas, \(2{,}47\) na Totalmente em GPU e \(2{,}93\) na CPU), nenhum par supera \(CD = 1{,}211\) (\figref{fig:cd-diagram-gap}) e os contrastes contra a CPU resultam em \(p_{\text{Holm}} = 1{,}00\). No Eixo~1 por instância (\tabref{tbl:axis1-wilcoxon} do \apref{apendice-resultados}), \(0\) das \(90\) comparações pareadas (\(6\) pares \(\times\) \(\totalinstancias\) instâncias) são significativas (\(p_{\text{Holm}} \ge 0{,}3418\)), com diferenças exatamente nulas entre as duas híbridas em todas as sementes (identidade bit a bit).
  \item
        \textbf{Tempo de execução:} A superioridade temporal das variantes aceleradas é confirmada em ambos os eixos. No Eixo~2 global, a ordenação foi idêntica nas \totalinstancias\ instâncias (\(\chi_F^2 = 45{,}00\), \(p < 0{,}001\); postos médios exatos \(1{,}00\), \(2{,}00\), \(3{,}00\) e \(4{,}00\)), todas as variantes em GPU superam a CPU via Wilcoxon com Holm--Bonferroni (\(p_{\text{Holm}} < 0{,}001\)) e o pós-teste de Nemenyi (\figref{fig:cd-diagram-time}, \(CD = 1{,}211\)) separa todos os pares distantes por duas ou mais posições (pares vizinhos ficam à distância \(1{,}00 < 1{,}211\)). No Eixo~1 por instância (\tabref{tbl:axis1-wilcoxon}), todos os \(6\) pares possíveis entre as \totalvariantes\ variantes diferem em todas as \totalinstancias\ instâncias (\(90/90\) testes com \(p_{\text{Holm}} \le 2{,}05 \times 10^{-7}\)), confirmando também a vantagem da Totalmente em GPU sobre a Híbrida Otimizada em cada instância individualmente.
\end{itemize}

\begin{figure}[htbp]
  \centering
  \caption{Diagrama de diferença crítica de Demšar (pós-teste de Nemenyi, \(CD = 1{,}211\), \(\alpha = 0{,}05\)) para o \textbf{tempo de execução} nas \totalinstancias\ instâncias: ordenação idêntica em todas elas (\(\chi_F^2 = 45{,}00\), \(p < 0{,}001\))}\label{fig:cd-diagram-time}
  \begin{adjustbox}{max width=\linewidth}
    % Auto-generated Demšar Critical Difference Diagram in TikZ
% Metric: Tempo de Execução
% Generated for k = 4 algorithms across N = 15 instances (CD = 1.211, alpha = 0.05, F_F = inf, p = 0)
\begin{tikzpicture}[
    scale=2.2,
    axis/.style={thick, -{Stealth[length=2.5mm]}},
    tick/.style={thick},
    subtick/.style={thin, gray},
    clique/.style={line width=2.4pt, color=blue!75!black, line cap=round},
    cdbar/.style={line width=1.5pt, {Bar[width=4mm]}-{Bar[width=4mm]}},
    algomarker/.style={circle, fill=black, inner sep=1.2pt},
    labelnode/.style={font=\small\bfseries, align=center}
]

    % 1. CD Reference Bar at the Top
    \draw[cdbar] (1.0, 1.10) -- node[above=2pt, font=\footnotesize\bfseries] {CD = 1,211 ($\alpha = 0{,}05$)} (2.211, 1.10);

    % 2. Main Horizontal Rank Axis (1.0 to k.0)
    \draw[axis] (0.8, 0.0) -- (4.3, 0.0) node[right, font=\footnotesize\bfseries] {Posto médio};

    % Major Ticks
    \draw[tick] (1.0, 0.08) -- (1.0, -0.08);
    \draw[tick] (2.0, 0.08) -- (2.0, -0.08);
    \draw[tick] (3.0, 0.08) -- (3.0, -0.08);
    \draw[tick] (4.0, 0.08) -- (4.0, -0.08);

    % Minor Sub-ticks (every 0.5)
    \draw[subtick] (1.5, 0.04) -- (1.5, -0.04);
    \draw[subtick] (2.5, 0.04) -- (2.5, -0.04);
    \draw[subtick] (3.5, 0.04) -- (3.5, -0.04);

    % 3. Algorithm Data Points and Drop Lines (staggered Demšar elbows avoid overlap on ties)
    \coordinate (A_FullGPUModular) at (1.000, 0.0);
    \node[algomarker] at (A_FullGPUModular) {};
    \draw[thick] (A_FullGPUModular) -- (1.000, -0.48) -- (0.72, -0.48);
    \node[labelnode, anchor=east] at (0.68, -0.48) {Totalmente em GPU\\(1,00)};
    \coordinate (A_HybridOptimized) at (2.000, 0.0);
    \node[algomarker] at (A_HybridOptimized) {};
    \draw[thick] (A_HybridOptimized) -- (2.000, -0.98) -- (0.72, -0.98);
    \node[labelnode, anchor=east] at (0.68, -0.98) {Híbrida Otimizada\\(2,00)};
    \coordinate (A_HybridNaive) at (3.000, 0.0);
    \node[algomarker] at (A_HybridNaive) {};
    \draw[thick] (A_HybridNaive) -- (3.000, -0.98) -- (4.28, -0.98);
    \node[labelnode, anchor=west] at (4.32, -0.98) {Híbrida Ingênua\\(3,00)};
    \coordinate (A_CPU) at (4.000, 0.0);
    \node[algomarker] at (A_CPU) {};
    \draw[thick] (A_CPU) -- (4.000, -0.48) -- (4.28, -0.48);
    \node[labelnode, anchor=west] at (4.32, -0.48) {CPU\\(4,00)};

    % Major Tick Labels (drawn after drop lines so white fill masks crossings)
    \node[below=2pt, font=\footnotesize, fill=white, inner sep=1.5pt] at (1.0, -0.08) {1,0};
    \node[below=2pt, font=\footnotesize, fill=white, inner sep=1.5pt] at (2.0, -0.08) {2,0};
    \node[below=2pt, font=\footnotesize, fill=white, inner sep=1.5pt] at (3.0, -0.08) {3,0};
    \node[below=2pt, font=\footnotesize, fill=white, inner sep=1.5pt] at (4.0, -0.08) {4,0};

    % 4. Non-Significant Cliques (|R_i - R_j| <= CD)
    \draw[clique] (1.000, 0.35) -- (2.000, 0.35);
    \draw[thin, blue!75!black, dashed] (1.000, 0.0) -- (1.000, 0.35);
    \draw[thin, blue!75!black, dashed] (2.000, 0.0) -- (2.000, 0.35);
    \draw[clique] (2.000, 0.50) -- (3.000, 0.50);
    \draw[thin, blue!75!black, dashed] (2.000, 0.0) -- (2.000, 0.50);
    \draw[thin, blue!75!black, dashed] (3.000, 0.0) -- (3.000, 0.50);
    \draw[clique] (3.000, 0.65) -- (4.000, 0.65);
    \draw[thin, blue!75!black, dashed] (3.000, 0.0) -- (3.000, 0.65);
    \draw[thin, blue!75!black, dashed] (4.000, 0.0) -- (4.000, 0.65);
\end{tikzpicture}

  \end{adjustbox}
\end{figure}

\begin{figure}[htbp]
  \centering
  \caption{Diagrama de diferença crítica de Demšar (pós-teste de Nemenyi, \(CD = 1{,}211\), \(\alpha = 0{,}05\)) para a \textbf{qualidade da solução (\emph{gap} percentual)}: nenhum par apresenta diferença estatisticamente significativa (\(\chi_F^2 = 2{,}69\), \(F_F(3, 42) = 0{,}8897\), \(p = 0{,}4543\))}\label{fig:cd-diagram-gap}
  \begin{adjustbox}{max width=\linewidth}
    % Auto-generated Demšar Critical Difference Diagram in TikZ
% Metric: Gap de Otimalidade
% Generated for k = 4 algorithms across N = 15 instances (CD = 1.211, alpha = 0.05, F_F = 0.890, p = 0.4543)
\begin{tikzpicture}[
    scale=2.2,
    axis/.style={thick, -{Stealth[length=2.5mm]}},
    tick/.style={thick},
    subtick/.style={thin, gray},
    clique/.style={line width=2.4pt, color=blue!75!black, line cap=round},
    cdbar/.style={line width=1.5pt, {Bar[width=4mm]}-{Bar[width=4mm]}},
    algomarker/.style={circle, fill=black, inner sep=1.2pt},
    labelnode/.style={font=\small\bfseries, align=center}
]

    % 1. CD Reference Bar at the Top
    \draw[cdbar] (1.0, 1.10) -- node[above=2pt, font=\footnotesize\bfseries] {CD = 1,211 ($\alpha = 0{,}05$)} (2.211, 1.10);

    % 2. Main Horizontal Rank Axis (1.0 to k.0)
    \draw[axis] (0.8, 0.0) -- (4.3, 0.0) node[right, font=\footnotesize\bfseries] {Posto médio};

    % Major Ticks
    \draw[tick] (1.0, 0.08) -- (1.0, -0.08);
    \draw[tick] (2.0, 0.08) -- (2.0, -0.08);
    \draw[tick] (3.0, 0.08) -- (3.0, -0.08);
    \draw[tick] (4.0, 0.08) -- (4.0, -0.08);

    % Minor Sub-ticks (every 0.5)
    \draw[subtick] (1.5, 0.04) -- (1.5, -0.04);
    \draw[subtick] (2.5, 0.04) -- (2.5, -0.04);
    \draw[subtick] (3.5, 0.04) -- (3.5, -0.04);

    % 3. Algorithm Data Points and Drop Lines (staggered Demšar elbows avoid overlap on ties)
    \coordinate (A_HybridNaive) at (2.300, 0.0);
    \node[algomarker] at (A_HybridNaive) {};
    \draw[thick] (A_HybridNaive) -- (2.300, -0.48) -- (0.72, -0.48);
    \node[labelnode, anchor=east] at (0.68, -0.48) {Híbrida Ingênua\\(2,30)};
    \coordinate (A_HybridOptimized) at (2.300, 0.0);
    \node[algomarker] at (A_HybridOptimized) {};
    \draw[thick] (A_HybridOptimized) -- (2.300, -0.98) -- (0.72, -0.98);
    \node[labelnode, anchor=east] at (0.68, -0.98) {Híbrida Otimizada\\(2,30)};
    \coordinate (A_FullGPUModular) at (2.467, 0.0);
    \node[algomarker] at (A_FullGPUModular) {};
    \draw[thick] (A_FullGPUModular) -- (2.467, -0.98) -- (4.28, -0.98);
    \node[labelnode, anchor=west] at (4.32, -0.98) {Totalmente em GPU\\(2,47)};
    \coordinate (A_CPU) at (2.933, 0.0);
    \node[algomarker] at (A_CPU) {};
    \draw[thick] (A_CPU) -- (2.933, -0.48) -- (4.28, -0.48);
    \node[labelnode, anchor=west] at (4.32, -0.48) {CPU\\(2,93)};

    % Major Tick Labels (drawn after drop lines so white fill masks crossings)
    \node[below=2pt, font=\footnotesize, fill=white, inner sep=1.5pt] at (1.0, -0.08) {1,0};
    \node[below=2pt, font=\footnotesize, fill=white, inner sep=1.5pt] at (2.0, -0.08) {2,0};
    \node[below=2pt, font=\footnotesize, fill=white, inner sep=1.5pt] at (3.0, -0.08) {3,0};
    \node[below=2pt, font=\footnotesize, fill=white, inner sep=1.5pt] at (4.0, -0.08) {4,0};

    % 4. Non-Significant Cliques (|R_i - R_j| <= CD)
    \draw[clique] (2.300, 0.35) -- (2.933, 0.35);
    \draw[thin, blue!75!black, dashed] (2.300, 0.0) -- (2.300, 0.35);
    \draw[thin, blue!75!black, dashed] (2.933, 0.0) -- (2.933, 0.35);
\end{tikzpicture}

  \end{adjustbox}
\end{figure}

A \figref{fig:scaling-time-gap} mostra como o tempo total de execução até o critério de parada (em escala logarítmica, painel superior) e o \emph{gap} percentual (painel inferior) evoluem conforme o número de cidades \(n\) cresce até \(\maiorinstancia\).
Para cada instância, as curvas representam a mediana das \totalrepeticoes\ repetições e as faixas sombreadas indicam o intervalo onde ficam os \(50\%\) centrais das execuções (\(Q_{25}\text{--}Q_{75}\)), evitando possíveis distorções causadas por execuções fora do padrão.
O gráfico confirma visualmente os dois resultados principais: uma separação consistente em ordens de grandeza no tempo entre as variantes e curvas de \emph{gap} praticamente sobrepostas.

\begin{figure}[htbp]
  \centering
  \caption{Evolução do tempo total de execução até a parada (escala logarítmica, painel superior) e do \emph{gap} percentual em relação ao ótimo da TSPLIB (painel inferior) conforme o número de cidades \(n\), exibindo a mediana e o intervalo dos \(50\%\) centrais (\(Q_{25}\text{--}Q_{75}\)) nas \totalrepeticoes\ repetições}\label{fig:scaling-time-gap}
  \includegraphics[width=\linewidth]{../../../benchmark_scripts/experiments-and-examples/chapter4/results/figures/scaling_time_gap.pdf}
\end{figure}

Dois comportamentos pontuais na \figref{fig:scaling-time-gap} decorrem diretamente do critério de parada (\(\varepsilon_{\text{alvo}} = 1\%\)) e da dificuldade própria de cada instância:

\begin{itemize}
  \item
        \textbf{Queda de \emph{gap} em \texttt{berlin52} e tempos curtos abaixo de \(100\) cidades:} Enquanto \texttt{att48} e \texttt{eil51} encerram na maioria das vezes assim que cruzam a tolerância \(\varepsilon_{\text{alvo}} = 1\%\) (\(106\) e \(116\) de \(\totalexecinstancia\) execuções, respectivamente), em \texttt{berlin52} a busca atinge o ótimo exato (\(\text{gap} = 0\%\)) em \(95\) de \(\totalexecinstancia\) execuções (em menos de \(5\) gerações em média), derrubando a mediana para perto de zero. Nessas instâncias pequenas, a variante Totalmente em GPU conclui em menos de \(100\)\,ms (\(14{,}0\)\,ms em \texttt{att48} a \(42{,}7\)\,ms em \texttt{gr96}), faixa em que o custo fixo de lançamento e sincronização de \emph{kernels} tem peso proporcional maior.
  \item
        \textbf{Redução de tempo entre \texttt{att532} (\(n = 532\)) e \texttt{si535} (\(n = 535\)):} Embora cada geração individual de \texttt{si535} leve ligeiramente mais tempo que em \texttt{att532}, todas as \(\totalexecinstancia\) execuções de \texttt{si535} atingem \(\varepsilon_{\text{alvo}} = 1\%\) em cerca de \(63\) gerações, contra \(173\) a \(194\) gerações em \texttt{att532} (onde \(47\) de \(\totalexecinstancia\) execuções encerram por estagnação acima de \(1\%\)). Como o gráfico plota o tempo total até a parada, esse menor número de gerações reduz o tempo total em \(46{,}8\%\) na CPU e em cerca de \(63\%\) nas GPUs --- e, pela proximidade entre \(n = 532\) e \(n = 535\) no eixo linear contínuo, a transição entre os dois pontos aparece como um degrau quase vertical.
\end{itemize}

\section{Divisão do tempo por etapa e análise de gargalos}\label{sec-44-decomposicao-fases}

Para explicar as diferenças de desempenho entre as arquiteturas sem depender apenas do tempo total, a \tabref{tbl:phase-timings} mostra a participação percentual de cada uma das oito etapas do ciclo evolutivo (\secref{apendice-telemetria}). Os valores de todas as \totalinstancias\ instâncias constam na \tabref{tbl:phase-timings-all} do \apref{apendice-resultados}.

\begin{table}[htbp]
  \centering
  \caption{Divisão percentual do tempo médio de execução entre as oito etapas do ciclo evolutivo (instâncias selecionadas)}\label{tbl:phase-timings}
  \begin{tabeladensa}
    \begin{tabular}{@{}llrccccccccc@{}}
      \toprule
      Instância         & Algoritmo         & Tempo (ms) & Inic. & Sel. & Cruz. & Mut. & Aval. & BL   & Sobr. & Sinc. & Total \\
      \midrule
      \textbf{berlin52} & CPU               & 1594,1     & 0,1   & 0,0  & 3,2   & 0,0  & 0,0   & 96,5 & 0,0   & 0,0   & 100,0 \\
                        & Híbrida Ingênua   & 510,7      & 0,2   & 0,1  & 10,5  & 0,1  & 0,2   & 88,7 & 0,1   & 0,0   & 99,9  \\
                        & Híbrida Otimizada & 78,8       & 2,2   & 0,6  & 79,0  & 0,9  & 2,3   & 13,8 & 0,6   & 0,2   & 99,6  \\
                        & Totalmente em GPU & 21,2       & 8,8   & 8,0  & 8,7   & 9,9  & 5,4   & 8,7  & 13,0  & 36,0  & 98,4  \\
      \midrule
      \textbf{kroA100}  & CPU               & 4805,8     & 0,0   & 0,0  & 2,5   & 0,0  & 0,0   & 97,3 & 0,0   & 0,0   & 100,0 \\
                        & Híbrida Ingênua   & 1183,2     & 0,1   & 0,1  & 10,7  & 0,1  & 0,2   & 88,6 & 0,1   & 0,0   & 100,0 \\
                        & Híbrida Otimizada & 186,8      & 1,1   & 0,6  & 73,1  & 0,8  & 1,8   & 21,6 & 0,6   & 0,2   & 99,7  \\
                        & Totalmente em GPU & 50,7       & 3,6   & 5,6  & 6,1   & 7,3  & 3,6   & 6,0  & 7,7   & 59,0  & 99,0  \\
      \midrule
      \textbf{d198}     & CPU               & 14658,7    & 0,0   & 0,0  & 2,0   & 0,0  & 0,0   & 97,9 & 0,0   & 0,0   & 100,0 \\
                        & Híbrida Ingênua   & 3079,6     & 0,1   & 0,1  & 12,0  & 0,1  & 0,3   & 87,2 & 0,1   & 0,0   & 100,0 \\
                        & Híbrida Otimizada & 546,1      & 0,5   & 0,5  & 63,7  & 0,7  & 1,4   & 32,0 & 0,7   & 0,1   & 99,7  \\
                        & Totalmente em GPU & 165,4      & 1,1   & 2,7  & 3,1   & 3,8  & 1,2   & 3,3  & 4,1   & 80,1  & 99,5  \\
      \midrule
      \textbf{pr264}    & CPU               & 23616,8    & 0,0   & 0,0  & 1,8   & 0,0  & 0,1   & 98,1 & 0,0   & 0,0   & 100,0 \\
                        & Híbrida Ingênua   & 4813,3     & 0,1   & 0,1  & 14,4  & 0,2  & 0,4   & 84,6 & 0,1   & 0,0   & 100,0 \\
                        & Híbrida Otimizada & 1020,2     & 0,4   & 0,5  & 60,9  & 0,7  & 1,3   & 35,3 & 0,7   & 0,1   & 99,8  \\
                        & Totalmente em GPU & 337,7      & 0,7   & 2,1  & 2,5   & 3,0  & 0,9   & 2,6  & 3,2   & 84,7  & 99,7  \\
      \midrule
      \textbf{att532}   & CPU               & 158238,2   & 0,0   & 0,0  & 1,7   & 0,0  & 0,1   & 98,2 & 0,0   & 0,0   & 100,0 \\
                        & Híbrida Ingênua   & 29333,2    & 0,0   & 0,1  & 10,9  & 0,1  & 0,5   & 88,1 & 0,1   & 0,0   & 100,0 \\
                        & Híbrida Otimizada & 6667,0     & 0,1   & 0,5  & 52,5  & 0,6  & 1,0   & 44,4 & 0,7   & 0,1   & 99,8  \\
                        & Totalmente em GPU & 3393,2     & 0,1   & 1,5  & 1,5   & 1,8  & 0,5   & 1,4  & 2,0   & 91,1  & 99,8  \\
      \midrule
      \textbf{si535}    & CPU               & 84187,4    & 0,0   & 0,0  & 1,2   & 0,0  & 0,1   & 98,7 & 0,0   & 0,0   & 100,0 \\
                        & Híbrida Ingênua   & 10965,6    & 0,0   & 0,1  & 10,8  & 0,1  & 0,6   & 88,2 & 0,1   & 0,0   & 100,0 \\
                        & Híbrida Otimizada & 2477,2     & 0,3   & 0,5  & 51,1  & 0,5  & 1,0   & 45,7 & 0,6   & 0,1   & 99,8  \\
                        & Totalmente em GPU & 1234,0     & 0,4   & 1,7  & 1,6   & 1,9  & 0,6   & 1,7  & 2,1   & 89,7  & 99,7  \\
      \midrule
      \textbf{pr1002}   & CPU               & 682575,6   & 0,0   & 0,0  & 1,4   & 0,0  & 0,1   & 98,4 & 0,0   & 0,0   & 100,0 \\
                        & Híbrida Ingênua   & 115353,9   & 0,0   & 0,1  & 9,2   & 0,1  & 0,7   & 89,7 & 0,2   & 0,0   & 100,0 \\
                        & Híbrida Otimizada & 25548,3    & 0,0   & 0,5  & 40,7  & 0,4  & 0,8   & 56,6 & 0,8   & 0,1   & 99,9  \\
                        & Totalmente em GPU & 14027,6    & 0,1   & 0,8  & 0,9   & 1,1  & 0,3   & 0,7  & 1,2   & 94,8  & 99,8  \\
      \bottomrule
    \end{tabular}
  \end{tabeladensa}
  \notatabela{Valores das etapas e da coluna \textit{Total} expressos em \% do tempo médio de execução (\textit{Tempo}); \textit{Inic.}: inicialização, \textit{Sel.}: seleção, \textit{Cruz.}: cruzamento, \textit{Mut.}: mutação, \textit{Aval.}: avaliação, \textit{BL}: busca local \(2\)-opt, \textit{Sobr.}: sobrevivência, \textit{Sinc.}: sincronização CPU--GPU. Na variante Totalmente em GPU, as colunas \textit{Sel.}, \textit{Cruz.}, \textit{Mut.}, \textit{Aval.} e \textit{BL} medem apenas o tempo gasto pela CPU para enfileirar chamadas assíncronas, enquanto a execução efetiva na GPU acumula-se integralmente em \textit{Sinc.} (até \(94{,}8\%\) em \texttt{pr1002}); por ser uma medição via relógio de CPU, essa decomposição não separa internamente o tempo gasto na GPU entre busca local \(2\)-opt e cruzamento OX.
    Pequenas diferenças abaixo de \(100\%\) na coluna \textit{Total} correspondem ao tempo residual fora das etapas cronometradas (\secref{apendice-telemetria}).
    Execuções abaixo de \(100\)\,ms sofrem influência do custo fixo de lançamento de \emph{kernels} (\(\approx 0{,}5\text{--}1{,}0\)\,ms).}
\end{table}

A leitura da \tabref{tbl:phase-timings}, em conjunto com o volume de transferências e lançamentos medido na \tabref{tbl:transfers-launches} (e os dados completos do \apref{apendice-resultados}), evidencia a mudança de gargalo prevista pela \textbf{Lei de Amdahl} (Equação~\eqref{eq:amdahl}, \secref{sec-263-cpu-vs-gpu}):

\begin{itemize}
  \item
        \textbf{CPU e Híbrida Ingênua (domínio da busca local):} Na CPU, a busca local \(2\)-opt (\textit{BL}) consome quase todo o tempo de execução --- de \(96{,}3\%\) a \(98{,}7\%\) (média de \(97{,}7\%\)) ---, enquanto o cruzamento (\textit{Cruz.}) ocupa apenas de \(1{,}2\%\) a \(3{,}5\%\). Na Híbrida Ingênua, mesmo com o \(2\)-opt rodando em GPU, lançar um \emph{kernel} e transferir cada rota separadamente pelo barramento PCIe mantém a busca local entre \(84{,}6\%\) e \(89{,}7\%\) do tempo total (média de \(87{,}8\%\)), com o cruzamento entre \(9{,}2\%\) e \(14{,}4\%\).
  \item
        \textbf{Híbrida Otimizada (inversão de gargalo entre cruzamento e busca local):} Ao agrupar toda a população em um único lançamento de \emph{kernel} por geração (\tabref{tbl:transfers-launches}), o tempo da busca local cai drasticamente (de \(103\,408{,}6\)\,ms para \(14\,406{,}3\)\,ms em \texttt{pr1002}). Com isso, pela Lei de Amdahl, a parte que permaneceu na CPU passa a dominar: o cruzamento por ordem (\textit{Cruz.}) sobe para \(40{,}7\%\) (\texttt{pr1002}) a \(79{,}5\%\) (\texttt{eil51}) do tempo total, com média de \(62{,}2\%\) no conjunto. Mesmo vetorizado em CPU, o cruzamento continua sendo a etapa mais cara até \(n = 666\) (\(52{,}4\%\)); somente nas duas maiores instâncias (\texttt{rat783}, com \(51{,}3\%\) de \textit{BL} contra \(45{,}8\%\) de \textit{Cruz.}, e \texttt{pr1002}, com \(56{,}6\%\) contra \(40{,}7\%\)) a busca local volta a liderar.
  \item
        \textbf{Totalmente em GPU (execução assíncrona e significado da coluna \textit{Sinc.}):} Ao levar seleção, cruzamento, mutação e sobrevivência para a GPU, o gargalo de cruzamento em CPU desaparece. Isso explica o ganho adicional de \(5{,}20\times\) sobre a Híbrida Otimizada em \texttt{eil51} --- onde o cruzamento consumia \(79{,}5\%\) do tempo ---, que se reduz para \(1{,}82\times\) em \texttt{pr1002}, onde a busca local já representava \(56{,}6\%\), conforme a Equação~\eqref{eq:amdahl}.
        Na leitura da \tabref{tbl:phase-timings}, é importante notar que o disparo de \emph{kernels} CUDA é \textbf{assíncrono}: a CPU apenas envia a fila de operações para a GPU em microssegundos (por isso \textit{Sel.}, \textit{Cruz.}, \textit{Mut.}, \textit{Aval.} e \textit{BL} marcam de \(0{,}3\%\) a \(10{,}0\%\)) e só pausa ao final da geração (\textit{Sinc.}) para aguardar a placa concluir todos os cálculos e então checar o critério de parada.
        Assim, os \(36{,}0\%\) a \(94{,}8\%\) exibidos em \textit{Sinc.} \textbf{não} são tempo ocioso de espera: eles correspondem ao \textbf{tempo real de processamento de todos os \emph{kernels} executando na GPU} (busca local \(2\)-opt e operadores genéticos), agrupados nessa coluna porque o cronômetro roda no relógio da CPU sem forçar pausas de espera entre um operador e outro (\secref{apendice-telemetria}).
\end{itemize}

Para quantificar diretamente o tráfego CPU--GPU e a frequência de chamadas de \emph{kernel} por trás desses perfis temporais, a \tabref{tbl:transfers-launches} apresenta as médias por repetição nas sete instâncias selecionadas (os valores das \totalinstancias\ instâncias constam na \tabref{tbl:transfers-launches-all} do \apref{apendice-resultados}).

\begin{table}[htbp]
  \centering
  \caption{Volume médio de transferências pelo barramento PCIe e número médio de lançamentos de \emph{kernel} por repetição (instâncias selecionadas)}\label{tbl:transfers-launches}
  \begin{tabeladensa}
    \begin{tabular}{@{}llrrrr@{}}
      \toprule
      Instância         & Algoritmo         & CPU\(\to\)GPU (MB) & GPU\(\to\)CPU (MB) & Total PCIe (MB) & Lançamentos     \\
      \midrule
      \textbf{berlin52} & Híbrida Ingênua   & \(0{,}31\)         & \(0{,}26\)         & \(0{,}57\)      & \(12\,373\)     \\
                        & Híbrida Otimizada & \(0{,}60\)         & \(0{,}27\)         & \(0{,}87\)      & \(54\)          \\
                        & Totalmente em GPU & \(0{,}14\)         & \(0{,}00\)         & \(0{,}14\)      & \(67\)          \\
      \midrule
      \textbf{kroA100}  & Híbrida Ingênua   & \(1{,}49\)         & \(1{,}06\)         & \(2{,}55\)      & \(26\,624\)     \\
                        & Híbrida Otimizada & \(2{,}36\)         & \(1{,}09\)         & \(3{,}45\)      & \(115\)         \\
                        & Totalmente em GPU & \(0{,}35\)         & \(0{,}00\)         & \(0{,}35\)      & \(150\)         \\
      \midrule
      \textbf{d198}     & Híbrida Ingênua   & \(9{,}34\)         & \(5{,}26\)         & \(14{,}60\)     & \(66\,389\)     \\
                        & Híbrida Otimizada & \(11{,}21\)        & \(5{,}31\)         & \(16{,}52\)     & \(286\)         \\
                        & Totalmente em GPU & \(0{,}98\)         & \(0{,}00\)         & \(0{,}98\)      & \(362\)         \\
      \midrule
      \textbf{pr264}    & Híbrida Ingênua   & \(22{,}54\)        & \(11{,}08\)        & \(33{,}62\)     & \(104\,960\)    \\
                        & Híbrida Otimizada & \(23{,}30\)        & \(11{,}17\)        & \(34{,}47\)     & \(452\)         \\
                        & Totalmente em GPU & \(1{,}55\)         & \(0{,}00\)         & \(1{,}55\)      & \(532\)         \\
      \midrule
      \textbf{att532}   & Híbrida Ingênua   & \(297{,}18\)       & \(96{,}53\)        & \(393{,}72\)    & \(453\,632\)    \\
                        & Híbrida Otimizada & \(197{,}05\)       & \(96{,}90\)        & \(293{,}95\)    & \(1\,950\)      \\
                        & Totalmente em GPU & \(6{,}17\)         & \(0{,}00\)         & \(6{,}17\)      & \(2\,723\)      \\
      \midrule
      \textbf{si535}    & Híbrida Ingênua   & \(106{,}06\)       & \(34{,}31\)        & \(140{,}38\)    & \(160\,341\)    \\
                        & Híbrida Otimizada & \(72{,}65\)        & \(34{,}44\)        & \(107{,}09\)    & \(690\)         \\
                        & Totalmente em GPU & \(4{,}73\)         & \(0{,}00\)         & \(4{,}73\)      & \(883\)         \\
      \midrule
      \textbf{pr1002}   & Híbrida Ingênua   & \(2\,669{,}68\)    & \(543{,}26\)       & \(3\,212{,}94\) & \(1\,355\,435\) \\
                        & Híbrida Otimizada & \(1\,099{,}67\)    & \(544{,}34\)       & \(1\,644{,}02\) & \(5\,825\)      \\
                        & Totalmente em GPU & \(18{,}69\)        & \(0{,}00\)         & \(18{,}69\)     & \(6\,880\)      \\
      \bottomrule
    \end{tabular}
  \end{tabeladensa}
  \notatabela{Médias sobre as \totalrepeticoes\ repetições.
    A variante CPU não aparece na tabela por não realizar transferências PCIe nem lançamentos de \emph{kernel} (\(0{,}00\)\,MB e \(0\) lançamentos).
    Eventuais diferenças de \(0{,}01\)\,MB no total decorrem do arredondamento a duas casas decimais.
    Na variante Totalmente em GPU e na Híbrida Otimizada, o contador em software da telemetria soma o tamanho da matriz \(D\) compartilhada em VRAM (\(4\,n_{coords}^2\) bytes) em cada estratégia que a referencia além do contexto central (acréscimo contábil de \(8\,n_{coords}^2\) bytes, ou \(8{,}03\)\,MB em \texttt{pr1002}, sobre o tráfego físico real no PCIe).}
\end{table}

Os números da \tabref{tbl:transfers-launches} evidenciam três comportamentos estruturais:

\begin{enumerate}
  \item
        \textbf{Lançamentos e volume PCIe entre as duas híbridas:} A Híbrida Otimizada reduz os lançamentos em cerca de \(230\times\) frente à Híbrida Ingênua (de \(1\,355\,435\) para \(5\,825\) em \texttt{pr1002}). No volume \(\text{CPU} \to \text{GPU}\), a Ingênua reenvia \(D\) e as rotas uma vez por geração (proporcional a \(n_{coords}^2 + n_{pop}\,n_{coords}\)), enquanto a Otimizada mantém \(D\) fixa em VRAM, mas envia a matriz \(P(t)\) duas vezes por geração --- busca local e avaliação (proporcional a \(2\,n_{pop}\,n_{coords}\)).
        Por isso, para \(n_{coords} \le 264\) (próximo a \(n_{pop} = 256\)), a Otimizada trafega ligeiramente mais megabytes que a Ingênua --- sendo ainda assim mais de \(4{,}5\times\) mais rápida graças à queda de lançamentos ---, invertendo-se a partir de \texttt{lin318}, quando o termo quadrático \(n_{coords}^2\) faz a Ingênua disparar até \(2\,669{,}68\)\,MB em \texttt{pr1002} (contra \(1\,099{,}67\)\,MB da Otimizada).
  \item
        \textbf{Composição do tráfego na Totalmente em GPU:} A variante residente zera o retorno de rotas (\(0{,}00\)\,MB em \(\text{GPU} \to \text{CPU}\)) e derruba o tráfego físico real no PCIe em \texttt{pr1002} para \(10{,}66\)\,MB (\(301\times\) menor que a Híbrida Ingênua e \(154\times\) menor que a Híbrida Otimizada). O total de \(18{,}69\)\,MB reportado na tabela em \texttt{pr1002} soma três parcelas:
        \begin{itemize}
          \item \(5{,}12\)\,MB de carga inicial única (\(D\): \(4{,}02\)\,MB, \(P(0)\): \(1{,}03\)\,MB e \(N_k\): \(0{,}08\)\,MB);
          \item \(5{,}54\)\,MB acumulados de sorteios aleatórios enviados da CPU (\(\approx 11{,}26\)\,KB/geração nas \(491{,}4\) gerações médias, o que também explica \texttt{att532} superar \texttt{si535} em tráfego);
          \item \(8{,}03\)\,MB de recontagem contábil do ponteiro compartilhado de \(D\) pela telemetria (sem transferência física no PCIe, conforme a nota da \tabref{tbl:transfers-launches}).
        \end{itemize}
  \item
        \textbf{Lançamentos de \emph{kernel} na Totalmente em GPU:} O total é ligeiramente superior ao da Híbrida Otimizada (\(6\,880\) contra \(5\,825\) em \texttt{pr1002}) apenas porque seleção, cruzamento e mutação também rodam como \emph{kernels} na GPU, em vez de permanecerem na CPU.
\end{enumerate}

As implicações arquiteturais desses perfis de comunicação e memória são discutidas no \capref{sec:discussao}.
